
\label{sec:NPhard}
\subsection{Budgeted Matroid Independent Set}

\label{sec:bm_np}


\lemsatbm*

%\begin{proof}
%	We use a reduction from the SAT-matroid problem. Let $I \in \{0,1\}^*$ be a SAT-matroid instance; for simplicity, let $A = A_I, n = n(A), k = k_I$, and $\alpha = \alpha_I$.  
%	%
%	%If $k \notin \mathbb{N}_{>0}$, $\alpha \notin \left[ n^2 \right]$, or $A$ is not a SAT instance with $n$ variables, we can verify it in time $\io$ and decide correctly the SAT-matroid instance $I$ in time $|I|^{O(1)}$. Thus, for the following, assume that $n,k \in \mathbb{N}_{>0}, \alpha \in \left[ n^2 \right]$,  and that $A$ is a SAT instance with $n$ variables. %$I = \left(n,k,\alpha,\mm\right)$ be a $\Pi$-matroid instance. 
%	We define the $\fs$-BM instance $$U = (I,w_{n},w_{n},B_{n,k,\alpha}).$$ Note that $\fs (I) = M_{n,k,\alpha}(\mathcal{S}(A))$ by \Cref{def:SATencoder}. %is an encodable matroid by \Cref{lem:encodable}; 
%	Moreover, by \Cref{lem:encodable}, \eqref{eq:w}, and \eqref{eq:B} it follows that $U$ can be constructed from $I$ in time $|I|^{O(1)}$.   
%	%
%	%For all $S \subseteq [n]$ it holds that $w_{n}(S) \in \mathbb{N}$ and that $$\OPT \left( U\right) \leq n \cdot \max_{i \in [n]} p(i) = n^2 = \left| U \right|^{O(1)}.$$  
%	%
%	To show that $U$ is a scaled $\fs$-BM instance, observe that for all $S \subseteq [n]$ it holds that $w_{n}(S) \in \mathbb{N}$ and that $\OPT \left( U\right) = \left| U \right|^{O(1)}$ by \Cref{lem:OPTu}. Therefore, if $\fs$-BM has an FPTAS, by \Cref{lem:FPTAStoEXACT} there is an exact algorithm $\cA$ that finds in time $\left| U \right|^{O(1)} = \io$ an optimal solution for $U$. Clearly, $\OPT(U) \leq B_{n,k,\alpha}$ since the profit and the cost of each element are the same. Therefore, by \Cref{cor:iff} the algorithm $\cA$ on instance $U$ returns a solution for $U$ with profit $B_{n,k,\alpha}$ if and only if $I$ has a solution. %Hence, we can decide $I$ by computing $\cA$ on $U$ and return that $I$ has a solution if and only if $\cA$ returns that $U$ has a solution for $U$ with profit $B_{n,k,\alpha}$Thus, $\cA$ is a polynomial algorithm that decides the SAT-matroid problem, 
%	Hence, by computing $\cA$ on $U$ and getting a solution $S$, we can decide correctly $I$ by returning that $I$ has a solution if and only if $w_n(S) = B_{n,k,\alpha}$. The above algorithm decides the SAT-matroid problem, which is impossible by \Cref{cor:Strong}, unless P=NP. 
%\end{proof}



\begin{proof}
	
	
	
	%	\noindent{\bf Proof of \Cref{lem:hard}:}
	Assume that P$\neq$NP and assume towards a contradiction that there is a Fully PTAS $\cA$ for $\fs$-BM. We use $\cA$ to decide the $\Pi$-matroid problem.  Let $I \in \{0,1\}^*$ be a SAT-matroid instance. Consider the following algorithm $\cB$ that decides $I$. For simplicity, let $A = A_I, n = n(A), k = k_I$, and $\alpha = \alpha_I$.  %let $n = n_I,k = k_I,\alpha = \alpha_I$, and $A = A_I$. 
	\begin{enumerate}
		\item Construct the $f_{\textnormal{SAT}}$-BM instance $%U_{n,k,\alpha}(\Pi) 
		U= (I,w_{n},w_{n},B_{n,k,\alpha})$. Note that $\fs (I) = M_{n,k,\alpha}(\mathcal{S}(A))$ by \Cref{def:SATencoder}. %is an encodable matroid by \Cref{lem:encodable}; 
		
		
		%; note that the elements of $U$ are $[n] = \{1,\ldots, n \}$, the independent sets are the independent sets of the $\Pi$-matroid $\mm$ (i.e., $\cI_{n,k,\alpha}(\Pi)$), the standard weight function $w_n$ is both the profit and the cost function, and the budget is simply the standard budget with parameters $n,k,\alpha$. %Let $U = U_{n,k,\alpha}(\Pi)$ for simplicity.
		\item Define $\eps = \frac{1}{6 \cdot n^3}$ and compute $\cA$ on $U$ and $\eps$; let $S$ be the returned solution.  
		\item Return that $I$ has a solution if and only if $w_n(S) = B_{n,k,\alpha}$. 
	\end{enumerate}
	
By \Cref{lem:encodable}, \eqref{eq:w}, and \eqref{eq:B} it follows that $U$ can be constructed from $I$ in time $|I|^{O(1)}$.   In addition, %Note that $U$ can be constructed from $I$ in time $\io$ and that 
$\frac{1}{\eps} = n^{O(1)} =|I|^{O(1)}$; thus, the running time of $\cB$ on $I$ is $|I|^{O(1)}$, since $\cA$ is an FPTAS. We now show correctness. %For the correctness, %we show that $\cB$ returns that $I$ has a solution if and only if $I$ indeed has a solution. 
	\begin{itemize}
		\item If $\cB$ returns that $I$ has a solution then $w_n(S) = B_{n,k,\alpha}$; since $S$ is a solution for $U$ it also holds that $S \in \cI_{\fs (I)} = \cI_{n,k,\alpha}(\mathcal{S}(A))$. Thus, by \Cref{cor:iff} it holds that $I$ has a solution.  
		\item If $I$ has a solution then there is $T \in \mathcal{L}_{n,k,\alpha}(\mathcal{S}(A))$; thus, by  \Cref{cor:iff} it follows that $w_n(T) = B_{n,k,\alpha}$ implying that $T$ is a solution for $U$.  Therefore, %w.p. at least $\frac{1}{2}$: 
		\begin{equation}
			\label{eq:sOPT}
			B_{n,k,\alpha} \geq \OPT(U) \geq w_n(S) \geq (1-\eps) \cdot w_n(T) \geq  w_n(T)-\frac{1}{6 \cdot n^3} \cdot 3 \cdot n^3 = B_{n,k,\alpha}-\frac{1}{2}.
		\end{equation} 	The first inequality holds since %$w_n(T) = B_{n,k,\alpha}$. In addition, $\OPT(U) \leq B_{n,k,\alpha}$ because 
		the profit and the cost of each element are the same, and the cost of a solution cannot exceed the budget. The third inequality follows because $S,T$ are solutions for $U$, and $S$ is computed by an FPTAS with the parameter $\eps$. The fourth inequality follows from \Cref{lem:OPTu} and by the selection of $\eps$. Observe that the value of $S$ is an integer by \eqref{eq:w}; thus, by \eqref{eq:sOPT}, %w.p. at least $\frac{1}{2}$ 
		it holds that $w_n(S) = B_{n,k,\alpha}$. Then, $\cB$ returns that $I$ has a solution.  
	\end{itemize}
%	Hence, %by computing $\cA$ on $U$ and getting a solution $S$, 
%	the algorithm $\cB$ is an solves the SAT matroid problem correctly. %if and only if $w_n(S) = \OPT(U)$, which holds w.p. at least $\frac{1}{2}$. 
Hence $\cB$ is %by returning that $I$ has a solution if and only if $w_n(S) = B_{n,k,\alpha}$. Thus, we 
	a polynomial time algorithm for the  SAT-matroid problem, contradicting \Cref{cor:Strong}. 
\end{proof}



\subsection{Exact Weight Matroid Independent Set}

Similarly to $f$-BM, we can define a variant of EMI for every matroid encoder $f$.

	\problem{ {\sc $f$-Exact Matroid Independent Set} ($f$-EMI)}{
	{\bf Encoder} & $f:\{0,1\}^* \rightarrow 2^{\mathbb{N}} \times 2^{2^{\mathbb{N}}}$ is a matroid encoder. \\
	{\bf Instance} &  $U =  (I,w,k,T)$, where $I \in \{0,1\}^*$, $w:E_{f(I)} \rightarrow \mathbb{N}$ gives the elements weights, $k\in \mathbb{N}$ is  the target cardinality, and $T \in \mathbb{N}$ is the target weight. \\
	{\bf Solution }& $S\in \cI_{f(I)}$ such that $|S|=k$ and $w(S)=T$.\\
{\bf Objective} & Decide if there exists a solution.
}



\begin{lemma}
	\label{lem:f-EMI}
	Unless \textnormal{W[1]=FPT}, for any computable function $f:\mathbb{N} \rightarrow \mathbb{N}$ there is no algorithm that decides \textnormal{EMI} in time $f(k) \cdot \left( T\cdot |I|\right)^{O(1)}$, where $|I|$ is the encoding size of the instance, $k$ is the target cardinality and $T$ is the target weight.
\end{lemma}
\begin{proof}
	Assume that \textnormal{W[1]}$\neq$\textnormal{FPT}. 
	Assume towards a contradiction that the statement of the lemma does not hold. Thus, there is a computable  function $f:\mathbb{N} \rightarrow \mathbb{N}$, and an algorithm $\cA$ that decides $\fs$-\textnormal{EMI} in time $f(k) \cdot {(T \cdot |I|)}^{O(1)}$. We use $\cA$ to decide the SAT-matroid problem. 
Let $I \in \{0,1\}^*$ be a SAT-matroid instance. Consider the following algorithm $\cB$ that decides $I$. For simplicity, let $A = A_I, n = n(A), k = k_I$, and $\alpha = \alpha_I$. Assume that $k \leq n$ or the instance can be trivially decided. 
	\begin{enumerate}
		\item Construct the $\fs$-EMI instance $%U_{n,k,\alpha}(\Pi) 
		U= (I,w_{n},k,B_{n,k,\alpha})$. Note that $\fs (I) = M_{n,k,\alpha}(\mathcal{S}(A))$ by \Cref{def:SATencoder}. %Let $U = U_{n,k,\alpha}(\Pi)$ for simplicity.
		\item %Define $\eps = \frac{1}{6 \cdot n^3}$ and 
		Compute $\cA$ on $U$. %and let $S$ be the returned solution.  
		\item Return that $I$ has a solution if and only if $\cA$ returns that $U$ has a solution.   %$w_n(S) = B_{n,k,\alpha}$. 
	\end{enumerate}
	To show $\cB$ indeed solves the SAT-matroid instance, consider the following cases. 
	\begin{itemize}
		\item If  $\cB$ returns that $I$ has a solution then $w_n(S) = B_{n,k,\alpha}$ and $S \in \cI_{\fs (I)} = \cI_{n,k,\alpha}(\mathcal{S}(A))$. Thus, by \Cref{cor:iff} it holds that $I$ has a solution.  
		\item If $I$ has a solution then there is $T \in \mathcal{L}_{n,k,\alpha}(\mathcal{S}(A))$; thus, by  \Cref{cor:iff} it follows that $w_n(T) = B_{n,k,\alpha}$. Also, by the definition of $\mathcal{L}_{n,k,\alpha}(\mathcal{S}(A))$ it holds that $|T|=k$.  Therefore $T$ is a solution of $U$. Therefore, $\cA$ returns that $U$ has a solution; thus, $\cB$ returns that $I$ has a solution. 
	\end{itemize}
It follows that $\cB$ is an algorithm for the SAT-matroid problem.

%	Note that $U$ can be constructed naively from $I$ in time $n+n \cdot (2n^2+n)+2n \cdot k+n^2 \leq 6 \cdot n^3$ (recall that the weights are encoded in unary and that $k \leq n$); thus, the running time of $\cB$ on $I$ is bounded by $6 \cdot n^3+f(k) \cdot n^{q} = f(k) \cdot n^{O(1)}$. 
	Note that $B_{n,k,\alpha} = 2k \cdot \textsf{sum}(n) \leq 2k \cdot n^2 \leq 2 \cdot n^3$ since $k \leq n$.  The instance $U$ can be trivially constructed in time $|I|^{O(1)}$, thus the overall running time of $\cB$ is $f(k)\cdot (2\cdot n^3\cdot |I|^{O(1)})^{O(1)} \leq f(k)\cdot |I|^{O(1)}$. That is, $\cB$ is a parameterized algorithm for the SAT-matroid problem, in contradiction to \Cref{lem:FPT}. Thus, the statement of the lemma holds.
\end{proof}

























\subsection{Constrained minimum Basis of a Matroid}

We proceed to define a variant of CMB in which the matroid is encoded as part of the instance.
\problem{ {\sc $f$-Constrained minimum Basis of a Matroid} ($f$-CMB)}{
	{\bf Encoder} & $f:\{0,1\}^* \rightarrow 2^{\mathbb{N}} \times 2^{2^{\mathbb{N}}}$ is a matroid encoder. \\
	{\bf Instance} &  
	$U = (I,c,w,B)$, where $I \in \{0,1\}^*$,  $c:E_{f(I)} \rightarrow \mathbb{N}$ is a cost function, $w:E_{f(I)} \rightarrow \mathbb{N}$ is a weight function and $B \in \mathbb{N}$ is a budget.\\
	{\bf Solution}& A basis $S \in \cI_{f(I)}, |S| = \max_{S' \in \cI_{f(I)}} |S'|$ such that $c(S) \leq B$.\\
	{\bf Objective} & Find a solution $S$ of $U$ with minimum weight $w(S)$.
	%	{\bf Model } &The independent sets $\cI$ is only accessible via a membership oracle; the encoding cost of the instance $U$ is the encoding cost of $(E,c,p,B)$, the instance without $\cI$.
}

\begin{lemma}
	\label{lem:CMB-NP}
	Unless \textnormal{P=NP}, there is no \textnormal{Fully PTAS} for $f_{\textnormal{SAT}}$-\textnormal{CMB}. 
	%If there is an algorithm for \textnormal{BM} in $X(n)$   
\end{lemma}
	\begin{proof}
		Assume that P$\neq$NP and assume towards a contradiction that there is a Fully PTAS $\cA$ for $f_{\textnormal{SAT}}$-CMB. We use $\cA$ to decide the SAT-matroid problem.  Let $I \in \{0,1\}^*$ be a SAT-matroid instance. Consider the following algorithm $\cB$ that decides $I$. For simplicity, let $A = A_I, n = n(A), k = k_I$, and $\alpha = \alpha_I$. 
	\begin{enumerate}
		\item Construct the $f_{\textnormal{SAT}}$-CMB instance $%U_{n,k,\alpha}(\Pi) 
		U= (I,c_n,w_{n},L_{n,k,\alpha})$. Note that $\fs (I) = M_{n,k,\alpha}(\mathcal{S}(A))$ by \Cref{def:SATencoder}.  %along with the membership oracle of $M_{n,k,\alpha}(\Pi)$; note that the elements of $U$ are $[n] = \{1,\ldots, n \}$, the independent sets are the independent sets of the $\Pi$-matroid $\mm$ (i.e., $\cI_{n,k,\alpha}(\Pi)$), the cost is the standard cost function, the weight is the standard weight function $w_n$, and the bound is the standard bound with parameters $n,k,\alpha$. %Let $U = U_{n,k,\alpha}(\Pi)$ for simplicity.
		\item Define $\eps = \frac{1}{6 \cdot n^3}$ and compute $\cA$ on $U$ and $\eps$; let $S$ be the returned solution.  
		\item Return that $I$ has a solution if and only if $w_n(S) = B_{n,k,\alpha}$. 
	\end{enumerate}
By \Cref{lem:encodable}, \eqref{eq:w}, and \eqref{eq:B} it follows that $U$ can be constructed from $I$ in time $|I|^{O(1)}$.   In addition, %Note that $U$ can be constructed from $I$ in time $\io$ and that 
$\frac{1}{\eps} =n^{O(1)} = |I|^{O(1)}$; thus, the running time of $\cB$ on $I$ is $|I|^{O(1)}$, since $\cA$ is an FPTAS. We now show correctness.
	%Note that $U$ can be constructed from $I$ in time $\io$ and that $\frac{1}{\eps} = |I|^{O(1)}$; thus, the running time of $\cB$ on $I$ is $|I|^{O(1)}$, since $\cA$ is a (randomized) FPTAS. We now show correctness. %For the correctness, %we show that $\cB$ returns that $I$ has a solution if and only if $I$ indeed has a solution. 
	\begin{itemize}
		\item If $\cB$ returns that $I$ has a solution. Then, $w_n(S) = B_{n,k,\alpha}$; since $S$ is a solution for $U$ it also holds that $S \in \cI_{\fs (I)} = \cI_{n,k,\alpha}(\mathcal{S}(A))$. Thus, by \Cref{cor:iff} it holds that $I$ has a solution.  
		\item If $I$ has a solution. Then, there is $T \in \mathcal{L}_{n,k,\alpha}(\mathcal{S}(A))$; thus, by  \Cref{cor:iff} it follows that $w_n(T) = B_{n,k,\alpha}$, which implies by \Cref{lemma:wl} that $c_n(T) \leq L_{n,k,\alpha}$. Then, since $T$ is a basis of $\mm$ (recall that $T \in \mathcal{L}_{n,k,\alpha}(\mathcal{S}(A))$), it follows that $T$ is a solution for $U$.  Hence,   
		\begin{equation}
			\label{eq:sOPTcmb}
			B_{n,k,\alpha} \leq \OPT(U) \leq w_n(S) \leq (1+\eps) \cdot w_n(T) \leq  w_n(T)+\frac{1}{6 \cdot n^3} \cdot 3 \cdot n^3 = B_{n,k,\alpha}+\frac{1}{2}.
		\end{equation} 	The first inequality follows from \Cref{lemma:wl} (note that $U$ has a solution since $T$ is a solution). %for every $S' \subseteq [n], |S'|<k$ it holds that $s$ every solution $S'$ for $U$ is a basis of $\mm$; thus, it follows that $w$ %$w_n(T) = B_{n,k,\alpha}$. In addition, $\OPT(U) \leq B_{n,k,\alpha}$ because 
		%	the profit and the cost of each element are the same, and the cost of a solution cannot exceed the budget. 
		The third inequality follows because $S,T$ are solutions for $U$, and $S$ is computed by an FPTAS with the parameter $\eps$. The fourth inequality follows from \Cref{lem:OPTu} and by the selection of $\eps$. Observe that the value of $S$ is an integer by \eqref{eq:w}; thus, by \eqref{eq:sOPTcmb}, %w.p. at least $\frac{1}{2}$ 
		it holds that $w_n(S) = B_{n,k,\alpha}$. Then, $\cB$ returns that $I$ has a solution.  
	\end{itemize}
	Hence, %by computing $\cA$ on $U$ and getting a solution $S$, 
	we decide $I$ correctly. %if and only if $w_n(S) = \OPT(U)$, which holds w.p. at least $\frac{1}{2}$. 
	It follows that $\cB$ is %by returning that $I$ has a solution if and only if $w_n(S) = B_{n,k,\alpha}$. Thus, we 
	a polynomial algorithm which decides the SAT-matroid problem, contradicting \Cref{cor:Strong}. 
\end{proof}























	\subsection{Knapsack Cover with a Matroid}

Finally, we define a variant of KCM with no oracle.

\problem{ {\sc $f$-Knapsack Cover with a Matroid} ($f$-KCM)}{
	{\bf Encoder} & $f:\{0,1\}^* \rightarrow 2^{\mathbb{N}} \times 2^{2^{\mathbb{N}}}$ is a matroid encoder. \\
	{\bf Instance} &  
	$U = (I,c,d,D)$, where $I \in \{0,1\}^*$,  $c:E_{f(I)} \rightarrow \mathbb{N}$ is a cost function, $d:E_{f(I)} \rightarrow \mathbb{N}$ is a size function, and $B \in \mathbb{N}$ is a demand.\\
	{\bf Solution}& $S \in \cI_{f(I)}$ such that $d(S) \geq D$.\\
	{\bf Objective} & Find a solution $S$ of $U$ with minimum cost $c(S)$.
	%	{\bf Model } &The independent sets $\cI$ is only accessible via a membership oracle; the encoding cost of the instance $U$ is the encoding cost of $(E,c,p,B)$, the instance without $\cI$.
}

\begin{lemma}
	\label{lem:KCM-NP}
	Unless \textnormal{P=NP}, there is no \textnormal{Fully PTAS} for $f_{\textnormal{SAT}}$-\textnormal{KCM}. 
	%If there is an algorithm for \textnormal{BM} in $X(n)$   
\end{lemma}

\begin{proof}
	Assume that P$\neq$NP and assume towards a contradiction that there is an FPTAS $\cA$ for $f_{\textnormal{SAT}}$-KCM. We use $\cA$ to decide the SAT-matroid problem. Let $I \in \{0,1\}^*$ be a SAT-matroid instance. Consider the following algorithm $\cB$ that decides $I$. For simplicity, let $A = A_I, n = n(A), k = k_I$, and $\alpha = \alpha_I$. 
	\begin{enumerate}
		\item Construct the $f_{\textnormal{SAT}}$-KCM instance $%U_{n,k,\alpha}(\Pi) 
	U= (I,w_n,w_{n},B_{n,k,\alpha})$. Note that $\fs(I) = M_{n,k,\alpha}(\mathcal{S}(A))$ by \Cref{def:SATencoder}. 
		\item Define $\eps = \frac{1}{6 \cdot n^3}$ and compute $\cA$ on $U$ and $\eps$; let $S$ be the returned solution.  
		\item Return that $I$ has a solution if and only if $w_n(S) = B_{n,k,\alpha}$. 
	\end{enumerate}
By \Cref{lem:encodable}, \eqref{eq:w}, and \eqref{eq:B} it follows that $U$ can be constructed from $I$ in time $|I|^{O(1)}$.   In addition, %Note that $U$ can be constructed from $I$ in time $\io$ and that 
$\frac{1}{\eps} =n^{O(1)} =  |I|^{O(1)}$; thus, the running time of $\cB$ on $I$ is $|I|^{O(1)}$, since $\cA$ is an FPTAS. We now show correctness.
	\begin{itemize}
		\item If $\cB$ returns that $I$ has a solution. Then, $w_n(S) = B_{n,k,\alpha}$; since $S$ is a solution for $U$ it also holds that $S \in \cI_{\fs (I)} =  \cI_{n,k,\alpha}(\mathcal{S}(A))$. Thus, by \Cref{cor:iff} it holds that $I$ has a solution.  
		\item If $I$ has a solution. Then, there is $T \in \mathcal{L}_{n,k,\alpha}(\mathcal{S}(A))$; thus, by  \Cref{cor:iff} it follows that $w_n(T) = B_{n,k,\alpha}$ implying that $T$ is a solution for $U$.  Therefore, 
		\begin{equation}
			\label{eq:sOPT}
			B_{n,k,\alpha} \leq \OPT(U) \leq w_n(S) \leq (1+\eps) \cdot w_n(T) \leq  w_n(T)+\frac{1}{6 \cdot n^3} \cdot 3 \cdot n^3 = B_{n,k,\alpha}+\frac{1}{2}.
		\end{equation} 	The first inequality holds since %$w_n(T) = B_{n,k,\alpha}$. In addition, $\OPT(U) \leq B_{n,k,\alpha}$ because 
		the size and the cost of each element are the same, and the size of a solution must be higher or equal from the demand. The third inequality follows because $S,T$ are solutions for $U$, and $S$ is computed by an FPTAS with the parameter $\eps$. The fourth inequality follows from \Cref{lem:OPTu} and by the selection of $\eps$. Observe that the value of $S$ is an integer by \eqref{eq:w}; thus, by \eqref{eq:sOPT}, %w.p. at least $\frac{1}{2}$ 
		it holds that $w_n(S) = B_{n,k,\alpha}$. Then, $\cB$ returns that $I$ has a solution.  
	\end{itemize}
	Hence, %by computing $\cA$ on $U$ and getting a solution $S$, 
	we decide $I$ correctly. %if and only if $w_n(S) = \OPT(U)$, which holds w.p. at least $\frac{1}{2}$. 
	It follows that $\cB$ is %by returning that $I$ has a solution if and only if $w_n(S) = B_{n,k,\alpha}$. Thus, we 
	a polynomial algorithm which decides the SAT-matroid problem, contradicting \Cref{cor:Strong}. 
\end{proof}















